% !TEX root = main.tex

\section{Complexity of choice functions}

In this section, we introduce the multifunctions, and we expose the pivotal result of this paper, that is seen to be a generalization of \cite[Theorem 2]{staiger2002kolmogorov}. We start biving the definition of a multifunction, and of a computable multivalued function.

\begin{definition}[Multivalued function]\label{def:multivalued_function}
    Let $X,Y$ be two sets. We write $f: X \rightrightarrows Y$ to denote the fact that $f$ is a multivalued function from $X$ to $Y$, i.e., for each $x \in X$, $f(x)$ is a subset of $Y$.
\end{definition}

\begin{definition}[Computable multivalued function]\label{def:computable_multivalued_function}
    Let $f: \{0,1\}^\ast \rightrightarrows \{0,1\}^\ast$ be a multivalued function. We say that $f$ is computable if there exists a computable function $g : \{0,1\}^\ast \to \{0,1\}^\ast$ such that 
    \begin{equation}
        g(x) = \langle  f(x) \rangle. 
    \end{equation}
\end{definition}

We now state the result on the algorithmic complexity of the multifunctions, which is a reformulation of classical results of nonprobabilistic statistics \cite[Section 5.5]{li2008introduction}.

\begin{theorem}[Computable multivalued functions and complexity]\label{thm:computable multivalued functions and complexity}
    Let $f: \{0,1\}^\ast \rightrightarrows \{0,1\}^\ast$ be a computable multivalued function and $x \in \{0,1\}^\ast$. For all $y \in f(x)$,
    \begin{equation}\label{eq:main:thm:computable multivalued functions and complexity}
        K(y) \leq K(x) + \log_2 \#f(x) + 2\log \log \#f(x) + \bigOw[]{x}.
    \end{equation}
\end{theorem}
\begin{proof}
    Let $f: \{0,1\}^\ast \rightrightarrows \{0,1\}^\ast$ be a computable multivalued function. Then, by Definition \ref{def:computable_multivalued_function}, there exists $g : \{0,1\}^\ast \to \{0,1\}^\ast$ computable, such that 
    \begin{equation}
        g(x) = \langle  f(x) \rangle,
    \end{equation}
    where $\bar y$ denotes the prefix-free encoding that we entroduced in the notation section. For $x \in \{0,1\}^\ast$, we let $y_1(x), y_2(x), \ldots, y_{\#f(x)}$ be the list of all the elements of $f(x)$, such that A REVOIR
    \begin{equation}
        g(x) = \coprod_{i=1}^{\#f(x)} \bar{y}_i(x),
    \end{equation}
    We can construct a computable function $h: \{0,1\}^\ast \to \{0,1\}^\ast$ such that
    \begin{equation}
        h(\langle x, i \rangle) = y_i(x),
    \end{equation}
    as follows. On input $\langle x, i \rangle$, use the function $g$ to generate the sequence $y = \coprod_{i=1}^{\#f(x)} \bar{y}_i(x)$. Then, use the sequence $y$ to enumerate the elements of $f(x)$ until reaching $\bar{y}_i(x)$, and output $y_i(x)$. Hence, since $h$ is computable, one has
    \begin{align}
        K[y_i(x)] = K[h(\langle x, i \rangle)] &\osetlem{\ref{lem:computable functions decrease complexity}}{\leq} K[\langle x, i \rangle] + \bigOw[]{x}\\
        &\overset{(a)}{\leq} K[x] + K[i] + 2\log K[i] + \bigOw[]{x},\\
        &\overset{(b)}{\leq} K[x] + \len i + 2 \log \len i + \bigOw[]{x}\\
        &\leq K[x] + \log i + 2 \log \log i + \bigOw[]{x},\\
        &\leq K[x] + \log \#f(x) + 2 \log \log \#f(x) + \bigOw[]{x},
    \end{align}
    for all $x \in \{0,1\}^\ast$ and $1 \leq i \leq \#f(x)$, where (a) and (b) follows from \cite{li2008introduction}, Example 2.1.5 and Theorem 2.1.2, respectively.
\end{proof}

As mentionned in previous section, we will consider relative algorithmic complexity defined in Definition \ref{def:relative kolmogorov complexity}. We further establish a result for multivalued functions that are computable relatively to a certain sequence $\xf \in \{0,1\}^\N$.
\begin{definition}(Relatively computable function)\label{def:relatively computable function}
    A function $g : \{0,1\}^\ast \to \{0,1\}^\ast$ is said to be computable relatively to $\yf\in \{0,1\}^\N$ if there exists an effective algorithm $h: \{0,1\}^\ast \to \{0,1\}^\ast$, such that for all $x \in \{0,1\}^\ast$, there exists a prefix $y$ of $\yf$ satisfying
    \begin{equation}\label{eq:main:def:relatively computable function}
        g(x) = h(\langle x, y\rangle).
    \end{equation}
    For $s \in \R$, we say that a multivalued function $f$ is computable relatively to $s$ if it is computable relatively to the greedy binary expansion $\yf$ of $\{s\} := s - \lfloor s \rfloor$.
\end{definition}
 Based on the definition of relatively computable function, we define relatively computable multivalued functions.
\begin{definition}[Relatively computable multivalued function]\label{def:relatively computable multivalued function}
    Let $f: \{0,1\}^\ast \rightrightarrows \{0,1\}^\ast$ be a multivalued function. We say that $f$ is computable relatively to $\yf\in \{0,1\}^\N$ if there exists a function $g : \{0,1\}^\ast \to \{0,1\}^\ast$ computable relatively to $\yf$ such that 
    \begin{equation}
        g(x) = \langle  f(x) \rangle. 
    \end{equation}
    For $s \in \R$, we say that a multivalued function $f$ is computable relatively to $s$ if it is computable relatively to the greedy binary expansion $\yf$ of $\{s\} := s - \lfloor s \rfloor$.
\end{definition}
We can finally derive a relation for relative algorithmic complexity of multivalued functions.
\begin{theorem}[Relatively computable multivalued functions and complexity]\label{thm:relatively computable multivalued functions and complexity}
    Let $f: \{0,1\}^\ast \rightrightarrows \{0,1\}^\ast$ be a multivalued function computable relatively to $\yf \in \{0,1\}^\N$ and $x \in \{0,1\}^\ast$. Then, 
    \begin{equation}\label{eq:main:thm:relatively functions computable multivalued  and complexity}
        K(y|\yf) \leq K(x|\yf) + \log_2 \#f(x) + 2\log \log \#f(x) + \bigOw[]{x},
    \end{equation}
    for all $x \in \{0,1\}^\ast$ and $y \in f(x)$. Moreover, (\ref{eq:main:thm:relatively functions computable multivalued  and complexity}) also holds if we replace $\yf \in \{0,1\}^\N$ by $\beta\in(1,2)$.
\end{theorem}
We omit the proof, as it is essentially the same as the proof of Theorem \ref{thm:computable multivalued functions and complexity}, except that we need to relativize every equation with respect to $\yf \in \{0,1\}^\N$.




% \begin{doubtfulsec}
    
%     \begin{definition}[Surrounding multivalued function]\label{def:surrounding_multivalued_function}
%         Let $X,Y$ be two sets. We say that a multivalued function $f: X \rightrightarrows Y$ surrounds a function $g : X \to Y$, if
%         \begin{equation}
%             g(x) \in f(x), \ \forall x \in X.
%         \end{equation}
%     \end{definition}
    
    
%     \begin{definition}\label{def:finite_multivalued_function}
%         Let $X,Y$ be two sets. We say that a multivalued function $f : X \rightrightarrows Y$ is finite if $f(x)$ is a finite set, for every $x \in X$.
%     \end{definition}
    
    
%     \begin{definition}[Ambiguous multivalued function]\label{def:ambiguous_multivalued_function}
%         Let $f: \{0,1\}^\ast \rightrightarrows \{0,1\}^\ast$ be a multivalued function and $\psi : \N \to \N$. We say that $f$ is $\psi$-ambiguous if 
%         \begin{equation}
%             \# f(x) \leq \psi(|x|), \ \forall x \in \{0,1\}^\ast.
%         \end{equation}
%     \end{definition}
    
    
%     \begin{theorem}[Ambiguity and complexity]
%         Let $f: \{0,1\}^\ast \rightrightarrows \{0,1\}^\ast$ be a $\psi$-ambiguous computable multivalued function and $x \in \{0,1\}^\ast$. For all $y \in f(x)$,
%         \begin{equation}
%             K(y) \leq K(x) + \log_2 \psi(|x|) + \bigOw{x}.
%         \end{equation}
%     \end{theorem}
    
%     \begin{corollary}
%         Let $g : \{0,1\}^\ast \to \{0,1\}^\ast$, and suppose that $g$ is surrounded by a $\psi$-ambiguous computable multivalued function. Then, for all $x \in \{0,1\}^\ast$,
%         \begin{equation}
%             K(g(x)) \leq K(x) + \log_2 \psi(|x|) + \bigOw{x}.
%         \end{equation}
%     \end{corollary}
%     \end{doubtfulsec}
    
%     \pdfcomment{Stay TO THE POINT: do not make too many undirect definitions}
    
%     \begin{doubtfulsec}
%     \begin{definition}
%         Let $b_1,b_2 \in (1,\infty)$. The greedy base conversion function is denoted by 
%         \begin{equation}
%             \begin{array}{cclc}
%                 C_{b_1,b_2} : & \subseteq \N^\N & \to & \N^\N\\
%                 & x & \mapsto &  \gf_{b_2}(\delta_{b_1}(x)).
%             \end{array}
%         \end{equation}
%     \end{definition}
%     \end{doubtfulsec}

